% GENERATED FILE. Edit canonical lessons, then run npm run generate:books.
% canonical-source-hash: fnv1a64:cbc0896ad4cd0986
% canonical-lessons: LA-C18-viridis-flavus

\chapter{Green and Yellow}
\label{ch:la-green-and-yellow}

This chapter is generated from the canonical micro-lessons used by Language Ladder.
Each section stays independently resumable and preserves the authored prerequisite order.

\section[viridis, flāvus]{viridis, flāvus — the kept word, and the word Romance abandoned twice over}

\label{lesson:LA-C18-viridis-flavus}



\begin{quote}
\textbf{Warm-up.} {[}PAUSE 2s{]} One of Latin's color words survived cleanly into Spanish. The other one is a genuine oddity: Latin had a perfectly good word for yellow — and virtually \textbf{no} Romance language kept it, though most of them agreed on a surprising replacement together.
\end{quote}

\begin{cousinweb}
\textbf{viridis} ("\textbf{green}") comes from \textbf{virēre}, "\textbf{to flourish, to be green, to be vigorous}" — and this is the \textbf{direct} source of Spanish's \emph{verde}, worn down by entirely ordinary sound change (syncope of the medial vowel: \emph{viridis} $\to$ \emph{virdis} $\to$ \emph{verde}). A clean, uneventful inheritance.
\end{cousinweb}

\begin{cousinweb}
\textbf{flāvus} ("\textbf{yellow, golden, blond}," often specifically of hair) is a genuine, well-attested Classical Latin color word, from \textbf{Proto-Indo-European} \emph{\textbf{*bhleh\textsubscript{3}w-}} ("shining, light-colored") — though \emph{lūteus} was Latin's more neutral, plain "yellow" word in everyday prose. Here's the honest oddity: almost \textbf{no} Romance language kept \emph{flāvus} as its everyday word for "yellow." Spanish and Portuguese independently reinvented \emph{amarillo}/\emph{amarelo} from \emph{amarus}, "bitter." But here's the genuinely surprising twist: \textbf{French}'s \emph{jaune}, \textbf{Italian}'s \emph{giallo}, and \textbf{Romanian}'s \emph{galben} all converge on a \textbf{different shared Latin word} instead — \textbf{galbinus}, which meant "\textbf{light green}," not yellow at all. So this isn't every daughter language scattering off in its own direction — it's closer to \textbf{two} major replacement stories: Ibero-Romance's shared "bitter" reinvention, and a much wider swath of Romance sharing \emph{galbinus} instead. \emph{Flāvus} itself didn't vanish completely — it survives today mainly in \textbf{Roman personal names}, like \textbf{Flāvius} and \textbf{Flāvia}.
\end{cousinweb}

\subsection*{Guided Practice}
{[}PAUSE 1s{]}

\begin{itemize}
\raggedright
  \item {[}YOU SAY: "viridis" — green, direct source of verde{]}
  \item {[}YOU SAY: "flāvus" — yellow, Latin's own word, almost universally abandoned{]}
  \item {[}YOU SAY: "galbinus" — "light green," yet somehow the shared source of French's jaune, Italian's giallo, AND Romanian's galben{]}
  \item {[}YOU SAY: "Flāvius" — a Roman name preserving flāvus, where the color word itself didn't survive{]}
\end{itemize}

\begin{tcolorbox}[breakable,colback=teal!4,colframe=teal!35!black,title={Wrap-up recall}]
{[}PAUSE 3s{]} What Latin verb gives \textbf{viridis}, and what Spanish word does it directly become? (\textbf{Virēre}, "to flourish" $\to$ \textbf{verde}.) Did most Romance languages keep \textbf{flāvus} as their word for yellow? (\textbf{No.}) What are the \textbf{two} main replacement stories, and which languages fall into each? (\textbf{Ibero-Romance} — Spanish/Portuguese — independently reinvented \emph{amarillo}/\emph{amarelo} from "bitter"; \textbf{French, Italian, and Romanian} all instead converge on the \emph{same} Latin word, \emph{galbinus}, "light green," not yellow at all.) Where does \emph{flāvus} mainly survive today? (\textbf{Roman personal names}, like \emph{Flāvius} and \emph{Flāvia}.)
\end{tcolorbox}
